% !TeX root = ../sections/02_exercises.tex
\subsection{2-A-5.}
\begin{exercise}
	数列
	\[
	\{a_n\}_{n\in\mathbb{N}}
	\]
	が Cauchy 列であることと，
	\[
	\lim_{N\to\infty}
	\sup_{n,m\geq N}|a_n-a_m|
	=0
	\]
	は同値であることを証明せよ。
\end{exercise}

\begin{background}
	数列 \(\{a_n\}\) が Cauchy 列であるとは，
	任意の \(\varepsilon>0\) に対して，
	ある \(N\in\mathbb{N}\) が存在して，
	任意の \(n,m\geq N\) に対して
	\[
	|a_n-a_m|<\varepsilon
	\]
	が成り立つことをいう。
	
	一方，
	\[
	\sup_{n,m\geq N}|a_n-a_m|
	\]
	は，\(N\) 番目以降の項どうしの距離の最大幅を表している。
	したがって，これが \(N\to\infty\) で \(0\) に近づくとは，
	十分後ろではすべての項どうしが一様に近いという意味である。
	
	これは Cauchy 列の定義と本質的に同じ内容である。
\end{background}

\begin{strategy}
	まず，
	\[
	S_N=\sup_{n,m\geq N}|a_n-a_m|
	\]
	とおく。
	示すべきことは，
	\[
	\{a_n\}\text{ が Cauchy}
	\quad\Longleftrightarrow\quad
	S_N\to 0
	\]
	である。
	
	まず，\(\{a_n\}\) が Cauchy 列であると仮定する。
	任意に \(\varepsilon>0\) を取る。
	Cauchy 条件より，ある \(N_0\in\mathbb{N}\) が存在して，
	任意の \(n,m\geq N_0\) に対して
	\[
	|a_n-a_m|<\varepsilon
	\]
	となる。
	したがって，\(N\geq N_0\) ならば
	\[
	S_N\leq \varepsilon
	\]
	である。
	よって
	\[
	S_N\to 0
	\]
	が従う。
	
	逆に，
	\[
	S_N\to 0
	\]
	を仮定する。
	任意に \(\varepsilon>0\) を取る。
	極限の定義より，ある \(N_0\in\mathbb{N}\) が存在して，
	\[
	N\geq N_0
	\quad\Longrightarrow\quad
	S_N<\varepsilon
	\]
	となる。
	
	特に \(N=N_0\) とすれば，
	\[
	\sup_{n,m\geq N_0}|a_n-a_m|<\varepsilon
	\]
	である。
	したがって，任意の \(n,m\geq N_0\) に対して
	\[
	|a_n-a_m|<\varepsilon
	\]
	が成り立つ。
	よって \(\{a_n\}\) は Cauchy 列である。
\end{strategy}